\section{Introduction}

Language models built with transformers \citep{devlin2018bert} have attained extensive success in natural language tasks such as language modeling \citep{radford2018improving}, text classification \citep{wang2018glue}, question answering \citep{rajpurkar2016squad}, and summarization \citep{liu2019fine}. The success is achieved by fine-tuning a big transformer model pre-trained with a large corpus. The target task for fine-tuning may only focus on a restricted scenario such as sentiment analysis \citep{socher2013recursive} and multiple-choice question inference \citep{zellers2018swag}. Having a big transformer model is often overkill for the target task and prohibits the model deployment to resource-constrained hardware. Therefore, language model compression raises immense interest.

The popular strategy creates a compact model from scratch \citep{jiao2019tinybert} or a subset of the big model's layers \citep{sun2019patient,sanh2019distilbert}, then pre-trains with a large corpus and distills knowledge from the big model. This process is called generic pre-training \citep{wang2020minilm,sun2019patient,sanh2019distilbert} and is necessary for a compact model to achieve good performance on the target tasks. However, the generic pre-training could still cost considerable computational resources. For example, it takes 384 NVIDIA V100 GPU hours to get the pre-trained TinyBERT \citep{jiao2019tinybert} on the Wiki corpus dataset. So it may not be affordable for everyone who wants to create a compact model. In contrast, another line of strategy, specifically low-rank factorization \citep{golub1971singular,noach2020compressing}, can potentially reduce a big model's parameters without the generic pre-training. Since the factorization aims to approximate the learned model parameters, the method has the nature of directly inheriting the knowledge of the big trained model. 

However, approximating the learned weights with standard factorization often loses most of the task performance. This work investigates this issue with the most popular strategy, which uses singular value decomposition (SVD) to compress the learned model weights. With SVD, the learned matrix is factorized into three matrices ($U$, $S$, $V$). The portion associated with small singular values will be truncated to produce a smaller version of factorized matrices. The multiplication of these smaller matrices will approximate the original one with fewer total parameters to achieve the model compression. In other words, SVD minimizes the reconstruction error with fewer parameters as its objective. However, this objective does not necessarily correlate to the ultimate goal of keeping task performance. Specifically, the SVD algorithm is biased to reconstruct the parameters associated with large singular values. As a result, the parameters mainly reconstructed by the ranks with small singular values will become the victim in the compression process. Are these victimized parameters less critical to achieving a good task performance? We argue that this is not true, and the optimization objective of SVD is not properly aligned with the target task objective. 
This paper is the first work to provide an empirical analysis of this issue, proposing a novel weighted SVD to mitigate it. 

Our weighted SVD addresses the above issue by assigning importance scores to the parameters. This score has to correlate to how much the task performance is affected by the parameter change. The Fisher information nicely fits into this purpose \citep{Pascanu14revisitingnatural}. Besides, the calculation of Fisher information is usually simplified to accumulating a parameter's squared gradient over the training dataset based on its task objective (\eg cross-entropy, regression error, etc.), conveniently providing the importance of each parameter in a model. Then we modify the optimization objective of factorization (\ie, reconstruction error) by multiplying it with Fisher information, providing a new objective that jointly considers matrix reconstruction error and the target task objective. 
% Figure \ref{fig:wsvd} illustrates the above concept, and we give a more formal description of the process in the later sections.

In summary, this work makes the following contributions: (1) we analyze the issue of mismatched objectives between factorization and the target task for model compression; (2) we propose a novel compression strategy with the SVD weighted by the Fisher information; (3) we perform extensive analysis on varied language tasks, showing our Fisher-weighted SVD can compress an already compact model, and it can achieve comparable compression rate and performance with methods that require an expensive generic model pre-training.